\section{Discussion and limitations}
\label{sec:discussion}
\paragraph{What SlowLab measures.} The three questions locate the models' weakness in the same place. They started below a standard grower setting and did not gain enough within the year to pass it (RQ1); they ran the four compartments as two synchronous batches and so never let in-season measurements change what they planted (RQ2); and the final reading step was not where they lost (RQ3). The bottleneck is how the agents used the facility as an instrument (where to start, what to try, when to plant), not how they summarised a year of outcomes.

\paragraph{What the baseline result means.} A simple trust-region search that starts from a standard grower setting and recommends only what it tried beat the models and reached most of the headroom found by an offline search. It brings domain knowledge the models lack, so the result shows how far general-purpose agents are from a domain-informed optimiser in this task, not that language models cannot do better; giving an agent the same starting point, or a scheduler that staggers experiments, are natural next tests on the same sites. Because the gap narrows when neighbouring spaces exchange more heat (Section~\ref{sec:robust}), the E1 numbers are not a general ranking of methods for greenhouse control.

\paragraph{Limitations.} SlowLab is a simulation study: the executor is not calibrated to a real greenhouse, its external check failed for temperature and humidity (Appendix~\ref{app:external}), and the sensor-noise model is a declared assumption. Results concern the tested facility, budget, action space and models. Three low-cost models were tested; stronger or reasoning-heavy models may behave differently, and GLM ran with mandatory capped reasoning. Only two environment sensitivities were run, so no robustness claim is made for other assumptions (boundary temperatures, heating and ventilation targets, sensor bias, control step). The baseline was redesigned after the pilot because the plain BO scored below the fixed reference; it was tuned only on development sites, but the redesign is a pilot-informed decision, and its radius was selected at the largest of the three preregistered candidates, so a wider search might make it stronger still. E3 drew each method's packets from different sites, so it bounds the effect of the reading step but cannot compare the quality of evidence across methods; a site-paired design would. The run was completed in two parts: the corrected baseline was rerun after the first part's results existed, and the corrected language-model branches were continued days after they started, so a provider-side change in between cannot be excluded (Appendix~\ref{app:original}).
